\documentclass[10pt,a4paper]{amsart}
\usepackage[margin=19mm]{geometry}
\usepackage{amsmath,amssymb,mathtools}
\usepackage[scr=rsfs,cal=euler]{mathalfa}
\usepackage{xfrac}
\usepackage[unicode,hidelinks,bookmarks=false]{hyperref}
\newcommand{\ie}{i.e.\ }
\newcommand{\cf}{cf.\ }
\newcommand{\resp}{respectively\ }
\newcommand{\Subsection}{\subsection}
\providecommand{\nobreakdash}{\nobreak\hbox{-}}
\setlength{\emergencystretch}{2em}
\multlinegap0pt
\usepackage{amssymb}
\usepackage{mathtools}
\newtheorem{lem}{Lemma}
\newtheorem{cor}{Corollary}
\newcommand{\B}{\mathbb{B}^n}
\newcommand{\R}{\mathbb{R}}
\renewcommand{\a}{\operatorname{a}}
\newcommand{\komI}[1]{{\bf \red [\refstepcounter{komcounter}\thekomcounter (by I): #1]}}
\newcommand{\red}{\color{red}}
\title{Hardy spaces and quasiregular mappings: averaged derivatives and the $\mathbb{BMO}$ case}
\author{Tomasz Adamowicz, Iv\'an Caama\~no}
\date{Source: arXiv:2605.13655v1 (CC BY 4.0)}
\begin{document}
\maketitle

\noindent\emph{Source and licence.} Hardy spaces and quasiregular mappings: averaged derivatives and the $\mathbb{BMO}$ case. Version arXiv:2605.13655v1 (CC BY 4.0), distributed by arXiv under the Creative Commons Attribution 4.0 International licence, http://creativecommons.org/licenses/by/4.0/, as recorded in the arXiv licence metadata for this exact version and shown on the abstract page https://arxiv.org/abs/2605.13655v1. Full licence notice: \texttt{licenses/arxiv-2605.13655-CC-BY-4.0.txt}.\par\smallskip

\noindent\textit{Editorial scope.} Counted unit: the article's cor cor:harnack-est (source0.tex lines 701--720). The article's complete original proof of it is delivered with it (source0.tex lines 721--756, one \texttt{proof} environment of 36 lines). No mathematics is written, repaired or re-derived: the statement and the proof are the article's own lines, in the article's own notation, and no retained line is substituted or re-flowed.  Retained uncounted as the source-local context the proof needs: lines 621--628; lines 623--626; lines 691--694.  Nothing was omitted from any retained span, so the package carries no omission record.  No retained line contains a citation, so the wrapper carries no bibliography and no bibliography entry.  No retained line contains a figure environment, an image-inclusion command or drawing code, so no image is delivered and the package declares no asset dependency.  Editorial formatting only, in addition: an AMS \texttt{amsart} preamble that restates the article's own declarations and macros used by the retained lines, namely the environments \texttt{lem}, \texttt{cor} on separate counters, together with the standard packages those lines need.  The retained environments print in this document as Lemma 1, Corollary 1 in order of appearance, whereas the article prints the same items with the numbering of its own full document.  The complete native source of the article as distributed in the arXiv e-print for this exact version is bundled unchanged as \texttt{sources/200526/source0.tex}.
\begin{lem}\label{lem:harnack}
	Let map $f:\B\rightarrow \R^n$ be such that its Jacoby matrix satisfies ${D\!f \in L^n_{loc}(\B)}$. Then for any $\lambda \in (0,1/2 )$ there exist constants $C_H=C_H(n, \lambda)>0$ and $0<\lambda_1<1<\lambda_2<2$, with $\lambda_1, \lambda_2$ depending only on $\lambda$, such that for all $x\in\B$ and $y\in\lambda B_x$
	\begin{align}
		& \frac{1}{C_H} \a_{f,\lambda_1}(y)\leq \a_f(x)\leq C_H\a_{f,\lambda_2}(y), \label{est1-harnack}\\
		& \frac{1}{C_H} \a_{f,\lambda_1}(x)\leq \a_f(y)\leq C_H\a_{f,\lambda_2}(x). \label{est2-harnack}
	\end{align}
Moreover, $\lambda_1$ and $\lambda_2$ in fact can be chosen so that $0<\lambda_1<1-\frac32\lambda$ and $\frac{2+2\lambda}{2-\lambda}<\lambda_2<2$.
\end{lem}

	\begin{align}
		& \frac{1}{C_H} \a_{f,\lambda_1}(y)\leq \a_f(x)\leq C_H\a_{f,\lambda_2}(y), \\
		& \frac{1}{C_H} \a_{f,\lambda_1}(x)\leq \a_f(y)\leq C_H\a_{f,\lambda_2}(x). 
	\end{align}

\begin{align}
 &\lambda_1:=\min \{ \lambda_1,\tilde{\lambda}_1\}<\min\left\{\frac{2-2\lambda}{2+\lambda}, \frac{2-3\lambda }{2}\right\}=\frac{2-3\lambda }{2}, \label{eq-lambda1}\\
 &\lambda_2:=\max\{ \lambda_2,\tilde{\lambda}_2\}>\max\left \{\frac{2+2\lambda}{2-\lambda}, \frac{3\lambda +2}{2}\right\}=\frac{2+2\lambda}{2-\lambda}. \label{eq-lambda2}
\end{align}

\begin{cor}\label{cor:harnack-est}
Let map $f:\B\rightarrow \R^n$ be such that its Jacoby matrix satisfies ${D\!f \in L^n_{loc}(\B)}$.
%Let $f:\B\rightarrow \R^n$ such that $Df$, the Jacoby matrix of $f$, exists almost everywhere in $\B$. 
Then the following Harnack estimates hold for the average derivatives $\a_{f,\lambda}$:
\begin{enumerate}
\item[(a)] For every $0<\lambda_1<2$ and $\lambda \in (0,  \lambda_1/3 )$ it holds for all $x\in \B , y\in\lambda B_x$  that
\begin{equation}
\a_{f,\lambda}(x)\leq \a_{f,\lambda_1}(y). 
\label{eq:Harnacklambda1prime}
\end{equation}
\item[(b)] For every $\lambda_2' \in (1,2)$ there exist parameters $\lambda\in (0,1/2)$ and $\lambda_2 <\lambda_2'$ such that for all $x\in \B$ and $y\in \lambda B_x$ it holds that
\begin{equation}
\a_f(y)\lesssim \a_{f,\lambda_2}(x) \lesssim \a_{f,\lambda_2'}(y). \label{cor12:est2}
\end{equation}
\item[(c)] For every $\lambda_1'\in (0,1)$ there exist parameters $\lambda\in (0,1/2)$ and $\lambda_1>\lambda_1'$ such that for all $x\in \B$ and $y\in \lambda B_x$ it holds that
\begin{equation}
 \a_{f,\lambda_1'}(y)\lesssim \a_{f,\lambda_1}(x) \lesssim \a_{f}(y). \label{cor12:est3}
\end{equation}
\end{enumerate}
\end{cor}

\begin{proof}
{\bf Part (a).} In order to prove assertion~\eqref{eq:Harnacklambda1prime} let us fix $\lambda_1 >0$. Let further $x\in \B$ and $\lambda , \lambda_1' \in (0,1)$ for $\lambda_1'$ depending on $x$ whose value will be determined below. Then, for $y\in \lambda B_x$ and $z\in \lambda_1' B_x$ we have
$$
|y-z|\leq |y-x|+|x-z|\leq \left( \frac\lambda 2 + \frac{\lambda_1^\prime}{2}\right) (1-|x|)\leq \left( \frac\lambda 2 + \frac{\lambda_1^\prime}{2}\right)\frac{2}{2-\lambda } (1-|y|),
$$
and the last inequality follows from the estimate $1-|x|\leq |x-y|+1-|y|\leq (1+\frac{\lambda}{2})(1-|y|)$.
Thus, in order to get that $\lambda_1' B_x\subset \lambda_1 B_y$ we need to ensure that
$$
\frac{\lambda + \lambda_1^\prime }{2-\lambda}\leq \frac{\lambda_1}{2}.
$$
This condition can be achieved, for instance, by setting $\lambda_1':=\lambda <\frac{\lambda_1}{3}$. Then $\lambda  B_x\subset \lambda_1 B_y$ and the assertion~\eqref{eq:Harnacklambda1prime} follows from the definition of $\a_{f,\lambda}$. 
\medskip
		
\noindent {\bf Part (b).} Let  $\lambda\in (0,1/2)$ and $\lambda_2\in (1,2)$ be as in Lemma~\ref{lem:harnack}; moreover, recall that $\lambda_2$ depends on $\lambda$. Then, for a fixed $x\in \B$ and all $y\in\lambda B_x$ and $z\in \lambda_2 B_x$ we have
	\begin{align*}
		|z-y|\leq |z-x|+|x-y|\leq \left( \frac{\lambda_2}{2}+\frac{\lambda}{2}\right) (1-|x|)\leq \left( \frac{\lambda_2}{2}+\frac{\lambda}{2}\right)\frac{2}{2-\lambda}(1-|y|).
	\end{align*}
In consequence, $z\in \lambda_2' B_y$ for every $\lambda_2' \in (1,2)$, provided that $\frac{\lambda_2+\lambda}{2-\lambda}\leq\frac{\lambda_2'}{2}$. Furthermore, since by~\eqref{eq-lambda2} we have that $\lambda_2> \frac{2+2\lambda}{2-\lambda}>1$, one may get any $\lambda_2' \in (1,2)$, upon choosing small $\lambda \approx 0$, resulting in big $\lambda_2\approx 1$.
% \komI{We actually have full range $\lambda_2^\prime \in (1,2)$ possible, so maybe the statement of the Lemma can be changed to something like for all $\lambda_2^\prime$. I don't know if this is also the case for $\lambda_1^\prime$.}
Therefore, we may choose $\lambda_2$ satisfying $\lambda_2<\lambda_2' \in (1,2)$ such that for all $x\in \B$ and $y\in \lambda B_x$ 
\begin{equation*}%\label{eq:Harnack2lem3.4} 
    \a_{f,\lambda_2}(x) \lesssim \a_{f,\lambda_2'}(y).
\end{equation*}
From this estimate, assertion~\eqref{cor12:est2} follows directly upon combining with the right-hand side estimate in~\eqref{est2-harnack}.
\medskip

\noindent {\bf Part (c).} Fix $\lambda_1' \in (0,1)$ and a point $x\in \B$. Then, for each $\lambda \in (0, \frac12)$ and $y\in \lambda B_x$ and $z\in \lambda_1'B_y$ the triangle inequality implies that
$$
|z-x|\leq \frac{\lambda_1^\prime}{2}(1-|y|)+\frac{\lambda }{2}(1-|x|)\leq  \left(\frac{\lambda_1'}{2}(1+\frac{\lambda}{2}) +\frac{\lambda}{2}\right) (1-|x|)
$$
Now let $\lambda_1$ be as in Lemma~\ref{lem:harnack}, then $z\in\lambda_1 B_x$ if
$$
\lambda_1^\prime (1+\frac{\lambda}{2} )+\lambda \leq \lambda_1.
$$
However, recall that by~\eqref{eq-lambda1} in Lemma~\ref{lem:harnack}, it holds that $\lambda_1<1-\frac32 \lambda$ and, thus, letting $\lambda \rightarrow 0^+$ allows for $\lambda_1\rightarrow 1^-$. By a limiting argument, since $\lambda_1^\prime <1$, we can choose $\lambda$ and $\lambda_1$ which verify both Lemma~\ref{lem:harnack} and the assertion~\eqref{cor12:est3}.
\end{proof}

\end{document}
